%% ============================================================
%%  Section: Main Results  (v2 of 20260404 — adds literature citations
%%  throughout interpretation; bank×quarter FEs in triple-diff;
%%  no triple-diff event-study figure)
%%  Depositor Discipline in Modern Banking (2026-04-04 draft)
%%
%%  Cross-references:
%%    tab:small_bank_did_time_deposits
%%    tab:small_bank_triple_diff_time_deposits
%%    tab:small_bank_did_financial_outcomes
%%    fig:es_time_deps_cecl
%%    fig:es_financial_outcomes_cecl
%%    fig:es_sophistication_split
%%    sec:identification (for design description)
%% ============================================================

\section{Main Results}
[@@@ section too long]
[@@@ avoid em dashes --- when possible]
[@@@ double check numbers cited here from tables. Results in latex/sections/results/tables_figures.tex]
\label{sec:results}

We now test whether the CECL Day-One disclosure triggered a disciplinary
response from uninsured depositors.  Our design isolates three empirical
predictions implied by the theory: (i) banks whose adoption-quarter CECL
adjustment revealed larger latent credit losses should experience relative
declines in uninsured deposit balances after the disclosure; (ii) the shift
should be concentrated in the uninsured portion of the deposit portfolio and not
replicated symmetrically in insured balances; and (iii) the informational cost of
losing uninsured funding should translate into higher funding costs and lower
profitability, the hallmark of a genuine disciplinary mechanism rather than a
mechanical accounting effect.  The results confirm all three predictions.  We
begin by establishing the quantity response in time deposits using a standard
difference-in-differences specification, then sharpen the test with a
triple-difference estimator that exploits within-bank-within-quarter variation to
absorb all bank-specific time-varying confounders, and finally document the
transmission to bank financial performance. [@@@ shorten]

\subsection{Depositor Quantity Response: Time Deposits}
\label{sec:results:time_deposits}
[@@@ no paragraph titles]
\paragraph{Why time deposits.}
We focus on time deposits as the primary measure of depositor discipline for
three reasons.  First, time deposits---certificates of deposit and other
fixed-term instruments---are held for investment rather than for transaction
purposes.  Unlike checking or savings balances, which are maintained partly for
liquidity convenience, a depositor who holds a time deposit has made a deliberate
choice to commit funds at a particular institution for a particular rate.
Discipline operates at the rollover decision: a depositor who declines to renew a
maturing CD, or who accepts a lower rate at a safer institution, is actively
expressing a risk judgment.  This is the mechanism that \citet{calomiris1991role}
formalize: demandable debt disciplines bank risk precisely because creditors can
withdraw at the rollover margin, not merely in the extreme event of imminent
failure.  The long literature on large CDs and bank risk-taking confirms that
deposit rates and quantities respond to risk signals at this margin
\citep{hannan1988bank, ellis1992does, park1995market}.[@@@ I am not sure about this last sentence, whether it is factually correct. double check]

Second, uninsured time deposits above the \$250{,}000 FDIC threshold are
disproportionately held by business depositors, municipalities, and financially
sophisticated households that have both the incentive and the capacity to monitor
bank risk.  These are precisely the depositors that theory identifies as the
relevant disciplinary agents: they face meaningful loss exposure in the event of
failure and are large enough to bear the fixed cost of acquiring and interpreting
financial disclosures \citep{flannery1998using}.  The CECL Day-One adjustment is
exactly the kind of structured, audited disclosure that such depositors can act
on \citep{gee2022decision, kim2026current}.

Third, time deposits are less subject to the payment-systems inertia that
attenuates the disciplinary response in demand accounts.  A business maintaining a
checking account at its primary bank faces switching costs---direct deposit
arrangements, bill-pay links, payroll systems---that make a rapid response to risk
signals costly even for a sophisticated depositor \citep{iyer2016tale}.  A
depositor holding a maturing CD faces essentially no switching cost: the
discipline mechanism operates at full force.  This is why prior evidence for
depositor quantity discipline is strongest in time-deposit data
\citep{maechler2006dynamic, egan2017deposit}. [@@@ these last three paragraphs are too long, don't be overconfident, we want to be confident, but not oversell or overstate]

\paragraph{Specification and controls.}
Table~\ref{tab:small_bank_did_time_deposits} presents difference-in-differences
estimates for uninsured and insured time deposits, separately, over the 97,064
bank-quarter observations in the community-bank panel.[@@@ sample size not necessary]  All specifications absorb
bank and calendar-quarter fixed effects.  Bank fixed effects remove all
time-invariant differences across institutions---balance-sheet composition,
business model, geographic market---so the estimates are identified from
within-bank changes over time.  Quarter fixed effects remove aggregate
macroeconomic forces common to all banks in a given quarter, including the Federal
Reserve's tightening cycle, which drove substantial deposit repricing across the
system after 2022 \citep{drechsler2017deposits}, and the post-pandemic
normalization of deposit flows. [@@@ avoid duplicating stuff from earlier in the paper. mention quick and short and move on if we have already talked about this stuff earlier in the paper. look at latex/main.tex and referenced tex files therein]

The specifications include three lagged controls.  Lagged log assets flexibly
captures changes in bank size that might independently affect deposit composition.
Lagged equity-to-assets controls for the bank's capital position in the prior
quarter.  This is particularly important because CECL itself mechanically reduced
regulatory capital for adopting banks through the retained-earnings charge---and
a weaker capital position could independently affect depositor confidence or bank
risk appetite \citep{bushman2015delayed}.  Including the lagged equity ratio
ensures that the CECL treatment coefficient captures the informational content of
the disclosure rather than the incidental capital reduction.\footnote{Results are
robust to excluding the equity control; the CECL coefficients are marginally
larger in magnitude, consistent with the equity ratio absorbing a small portion of
the disciplinary signal.}  Lagged interest expense scaled by assets controls for
the bank's ex-ante deposit pricing posture, ensuring that the post-adoption shift
in uninsured deposits is not confounded with differential starting points in
funding costs.  All controls enter with a one-quarter lag to avoid simultaneous
feedback.[@@@ duplicated informatoin from the previous section]

\paragraph{Uninsured time deposits.}
The interaction of the post-adoption indicator with the binary high-CECL treatment
(column~2) yields a coefficient of $-0.394$ (standard error $0.182$), significant
at the 5 percent level.  Relative to otherwise similar low-CECL banks, high-CECL
community banks saw their uninsured time deposits shrink by roughly 39 basis
points of total assets in the post-adoption period.  This magnitude is
economically meaningful. The median community bank in our sample holds
approximately \$500 million in total assets, so a 39-basis-point contraction in
uninsured time deposits implies roughly \$2 million in lost uninsured balances per
bank. [@@@ don't need this much info on economic magnitude. say something like it is x pct of sample mean]  Scaled to the top decile of the CECL distribution, where the average
adjustment exceeds 3.2 percent of pre-adoption equity, the implied balance
contraction represents a significant fraction of a typical community bank's
uninsured funding base---consistent with the literature's finding that even
moderate risk signals trigger measurable quantity responses among uninsured
depositors \citep{martinez2001depositors, maechler2006dynamic}. [@@@ too long]

The continuous-treatment estimate (column~1) is consistent: a
one-percentage-point increase in the CECL equity shock reduces uninsured time
deposits by $0.072$ percentage points of assets ($p < 0.05$).  At the top-decile
threshold the continuous specification predicts a contraction of approximately 23
basis points, reasonably close to the binary estimate of 39 basis points given
that the high-CECL indicator averages across banks just above the threshold. [@@@ this could be shorter. say consistent and robust]

\paragraph{Insured time deposits.}
The pattern for insured time deposits runs in the opposite direction.  Column~4
reports a high-CECL coefficient of $+0.712$ ($p < 0.05$), indicating that
high-exposure banks attracted substantially more insured time deposits after
adoption.  The contrast with the uninsured estimate is economically striking: the
same banks that shed 39 basis points of uninsured balances simultaneously added
71 basis points of insured balances.  The continuous insured estimate in column~3
is positive ($+0.070$) but does not reach conventional significance thresholds.

This asymmetric response is not incidental---it is the disciplinary mechanism in
operation.  As \citet{calomiris1991role} predict, when informed creditors
withdraw, banks must replace them with uninformed, protected depositors who lack
the incentive to monitor.  Uninsured depositors, exposed to loss in the event of
bank failure, reduced their balances upon observing that their bank carried
higher-than-expected latent credit losses.  High-CECL banks, facing the loss of
this funding, had to substitute toward insured deposits, bidding up rates on
certificates of deposit to attract replacement balances \citep{acharya2015crisis}.
The result is a compositional shift in the deposit portfolio toward the more
expensive, insurance-protected tier---precisely the dynamic that \citet{chen2022bank}
document when transparency about bank credit quality increases: informed,
uninsured depositors respond to newly revealed risk, and banks replace them with
insured funding at a premium.  We document the cost consequences of this
substitution in Section~\ref{sec:results:financial}. [@@@ don't need this much details. say insured response is muted which is what we would expect. sometimes it couuld go up mechanically, if depositors reduce the amount to fall below 250k]

\paragraph{Event-study dynamics and pre-trends.}
Figure~\ref{fig:es_time_deps_cecl} plots event-study coefficients for the window
of ten quarters on either side of CECL adoption ($k = -1$ omitted).  Two features
are diagnostic.  First, for both the binary (left panels) and continuous (right
panels) treatment, the pre-adoption coefficients for uninsured time deposits
cluster tightly around zero through all ten lead quarters and are jointly
statistically indistinguishable from the reference period.  High-CECL and low-CECL
banks were on essentially identical deposit trajectories before the disclosure---a
clean validation of the parallel-trends assumption.  Second, the post-adoption
break is sharp: the uninsured coefficient turns negative within the first adoption
quarter and remains persistently below zero, while the insured coefficient moves
in the opposite direction with comparable persistence.  The symmetry in timing and
the absence of anticipation effects reinforce the interpretation of the disclosure
as the operative trigger.[@@@ don't sound overconfident. like cleanly ]

\subsection{Triple-Difference: Within-Bank, Within-Quarter Composition} [@@@ don't make this a separate subsection. should be within the same section]
\label{sec:results:triple}

\paragraph{Design and fixed effects.}
The difference-in-differences estimates in Section~\ref{sec:results:time_deposits}
identify off within-bank changes over time, conditioning on quarter fixed effects
to absorb aggregate deposit trends.  A remaining concern is that high-CECL banks
may have faced bank-specific funding pressures after 2023---idiosyncratic to
their markets, liability structures, or local economic conditions---that shifted
the entire time-deposit portfolio independently of depositor discipline.
\citet{covitz2004market} note that empirical tests of market discipline must
carefully separate depositor responses to bank-specific information from aggregate
funding dynamics; if CECL simply depressed aggregate deposit demand at
high-exposure institutions, both insured and uninsured balances would move
together, and the DiD estimate would overstate the discipline effect.

To address this, we stack the uninsured and insured time-deposit series in a
single panel and estimate a triple difference.  The specification is: [@@@ break equation into two lines]
\begin{equation*}
  y_{idt}
  = \alpha_{it}
  + \beta_1 \cdot \mathrm{Uninsured}_d
  + \beta_2 \cdot \mathrm{Post}_{t} \times \mathrm{Uninsured}_d
  + \beta_3 \cdot \mathrm{CECL}_i \times \mathrm{Uninsured}_d
  + \beta_4 \cdot \mathrm{CECL}_i \times \mathrm{Post}_{t} \times \mathrm{Uninsured}_d
  + \varepsilon_{idt},
\end{equation*}
where $d \in \{\text{uninsured}, \text{insured}\}$ indexes the deposit type, and
$\alpha_{it}$ are \emph{bank $\times$ quarter} fixed effects.  The bank $\times$
quarter fixed effects absorb the full vector of bank-specific time-varying
confounders in each period: rate sensitivity, local economic conditions,
concurrent changes in the bank's risk appetite or liability structure, and any
other shock that affects the bank uniformly in a given quarter.  Identification
is therefore driven entirely by the differential behavior of uninsured versus
insured deposits \emph{within the same bank and the same quarter}.  The
triple-interaction coefficient $\hat{\beta}_4$ captures whether the within-bank
composition of uninsured relative to insured time deposits shifts for high-CECL
banks after adoption---a test that is immune by construction to any story in which
CECL depresses aggregate funding demand at high-exposure institutions.  This is
the cleanest test in the paper.[@@@ agian sounding overconfident. immune by construction to any story]

\paragraph{Results.}
Table~\ref{tab:small_bank_triple_diff_time_deposits} reports the estimates.  The
triple interaction---$\text{High CECL} \times \text{Post} \times
\text{Uninsured}$ in column~(2)---is $-1.252$ (standard error $0.340$),
significant at the 1 percent level.  For the continuous specification
(column~1), the analogous triple interaction is $-0.177$ (standard error
$0.058$), also significant at the 1 percent level.  Relative to the within-bank
composition in the pre-adoption period, and relative to low-CECL banks, high-CECL
banks experienced a 125-basis-point widening of the gap between their insured and
uninsured time-deposit balances as a share of assets.  

The $\text{Post} \times \text{Uninsured}$ coefficient in column~(2) is $+0.898$
($p < 0.01$), confirming that across the entire community-bank sample uninsured
time deposits rose after 2023Q1---plausibly reflecting the broader post-pandemic
rotation of deposits into higher-yielding time accounts as CD rates increased
\citep{drechsler2017deposits}.  Because bank $\times$ quarter fixed effects absorb
any level shift in total time-deposit demand at each bank in each quarter, the
triple-interaction estimate strips out this aggregate trend.  What remains is the
differential \emph{within-bank reallocation} away from uninsured deposits at the
banks whose credit-risk information was revised most adversely by CECL
adoption---the compositional signature of depositor discipline as distinct from
aggregate funding pressure \citep{chen2022bank}.[@@@ this paragraph is not clear, remove if not essential]

The continuous estimate reinforces the dose-response interpretation: a
one-percentage-point increase in the CECL equity shock widens the within-bank
insured-uninsured composition gap by 17.7 basis points of assets ($p < 0.01$),
confirming that the compositional shift is a smooth function of disclosure
severity rather than a threshold phenomenon concentrated at the extreme tail of
the CECL distribution.

\subsection{Financial Consequences: Funding Costs and Profitability} [@@@ shorten, don't sound overconfident]
\label{sec:results:financial}

Table~\ref{tab:small_bank_did_financial_outcomes} turns to the financial
consequences of the deposit reallocation.  The dependent variables are quarterly
interest expense scaled by assets (columns~1--2), return on assets
(columns~3--4), and net interest margin (columns~5--6).  All specifications
include bank and quarter fixed effects and lagged log assets and equity ratio.

\paragraph{Funding costs.}
Column~(2) shows that high-CECL banks paid $0.021$ percentage points more in
quarterly interest expense per dollar of assets after adoption ($p < 0.01$).
Annualized, this amounts to approximately 8--9 basis points of additional funding
cost relative to low-CECL peers.  The continuous estimate in column~(1) implies
that a one-percentage-point increase in the CECL equity shock raises quarterly
interest expense by 0.37 basis points of assets, consistent with the binary
estimate scaled at the top-decile threshold.

The mechanism is direct: having lost lower-cost uninsured deposits, high-CECL
banks replaced them with insured time deposits at rates sufficient to attract
depositors who would not otherwise bear uninsured credit risk at these
institutions.  This repricing is the price channel of depositor discipline---the
premium that banks pay for having their credit-risk composition revealed
\citep{nier2006market, billett1998cost}.  Prior tests of the price channel have
relied on subordinated debenture yields \citep{flannery1996evidence} or large
CD rates \citep{hannan1988bank, ellis1992does}; our estimates recover the same
signal through the interest expense line of the income statement, aggregating
across the full liability structure.

\paragraph{Profitability.}
The profitability effects confirm that the funding cost increase was not absorbed
by higher loan yields or cost-cutting elsewhere.  Return on assets fell by $0.038$
percentage points per quarter for high-CECL banks (column~4, $p < 0.01$), an
annualized drag of roughly 15 basis points.  For community banks operating at
typical ROA levels of 100--120 basis points annually in 2023--2024, this
represents a compression of approximately 12--15 percent of average earnings---a
penalty that is economically substantial and not a rounding error.  Net interest
margin also contracted by $0.024$ percentage points per quarter (column~6,
$p < 0.05$), an annualized effect of approximately 9 basis points.  The joint
decline in ROA and NIM confirms that the higher interest expense on the liability
side was not passed through to borrowers---consistent with a competitive loan
market in which community banks are price-takers on the asset side but must bid
up liabilities to compensate for the loss of uninsured depositors
\citep{drechsler2017deposits}.  The result is a persistent drag on profitability
that represents a genuine economic cost of carrying elevated credit risk, not
merely an accounting adjustment \citep{bushman2012accounting}.

The continuous treatment estimates scale proportionally: each percentage point of
CECL equity shock reduces quarterly ROA by $0.009$ percentage points ($p < 0.01$)
and quarterly NIM by $0.005$ percentage points ($p < 0.05$), reinforcing that the
effects are a smooth function of disclosure severity rather than a threshold
phenomenon.

\paragraph{Event-study dynamics.}
Figure~\ref{fig:es_financial_outcomes_cecl} plots event-study coefficients for
all three financial outcomes.  Pre-adoption leads are flat for interest expense,
ROA, and NIM across both treatment specifications, again supporting parallel
trends.  The post-adoption divergence for ROA and NIM is persistent through the
end of the sample horizon---consistent with a multi-year repricing dynamic in
which maturing insured CDs, raised to replace the outflowing uninsured balances,
continue to roll over at elevated rates in subsequent periods.  The interest
expense effect also persists, indicating that the cost of disciplinary funding
substitution is not merely a transitional adjustment but reflects a durable
repricing of community bank liabilities at high-CECL institutions.

\subsection{Heterogeneity by Depositor Sophistication} [@@@ shorten. at this late in the paper, readers don't want to spend a lot of time reading these sections.]
\label{sec:results:sophistication}

\paragraph{Motivation.}
The depositor discipline mechanism requires that depositors be both willing and
able to act on new risk information.  Not all depositors satisfy this condition
equally.  The theoretical and empirical literature emphasizes that discipline is
strongest when creditors are informed, face real loss exposure, and have credible
outside options \citep{flannery1998using, nier2006market}.  Depositors in
markets with higher financial sophistication---more college-educated households,
higher dividend income, greater financial market participation---have lower costs
of processing bank disclosures and stronger incentives to monitor
\citep{danisewicz2021debtholder}.  If the CECL-triggered response we document
reflects genuine depositor discipline rather than a mechanical or institutional
reaction, it should be more pronounced at banks whose depositors are better
equipped to act on financial information.  We test this prediction by splitting
the sample by a deposit-weighted measure of branch-market sophistication
constructed from the 2022 FDIC Summary of Deposits merged to ZIP-code-level
American Community Survey data \citep{ACS2022}.  Banks at or above the 75th
percentile of this measure are classified as operating in high-sophistication
markets.

\paragraph{Visual evidence: quantity response.}
Figure~\ref{fig:es_sophistication_split}, Panel~A overlays event-study
coefficients for uninsured time deposits separately for high- and
low-sophistication banks.  The pre-adoption leads for both groups are flat and
overlap closely, confirming common pre-trends even within the sophistication
split.  The divergence emerges after adoption and is markedly asymmetric: the
high-sophistication group drifts downward beginning around $k=+1$ and accelerates
steadily through the end of the sample horizon, reaching approximately $-1.5$ to
$-2$ percentage points of assets by quarters~$+9$ to~$+11$.  The
low-sophistication group remains clustered near zero across the entire post-period.
This pattern---a gradual, widening separation rather than a sharp initial
break---explains why the average triple-interaction coefficient $\text{Post}
\times \text{High CECL} \times \text{High Sophistication}$ is small and
statistically indistinguishable from zero ($-0.022$, SE~$= 0.468$): the full
post-period average pools early quarters in which the two groups have not yet
separated with later quarters in which the gap is large and economically
substantial.  The gradual dynamics are consistent with \citet{iyer2016tale}'s
finding that sophisticated depositors respond to solvency signals through
deliberate portfolio reallocation rather than an immediate panic---here, allowing
CDs to mature before redirecting funds rather than breaking term deposits at a
cost.

Panel~B tells a complementary story for insured time deposits: in
high-sophistication markets, high-CECL banks also attracted fewer insured deposits
than their low-sophistication counterparts ($-0.965$, SE~$= 0.677$,
$p = 0.154$).  While this coefficient falls short of conventional significance,
its direction is consistent with sophisticated depositors being less easily
recruited as replacement insured depositors---they demand more before committing
to insured CDs at institutions whose credit quality has been disclosed as weaker
than expected.

\paragraph{Price channel: funding costs and profitability.}
The sharpest statistical signal from the sophistication split appears in the
financial outcome panels.  Panel~C shows that high-CECL banks in
high-sophistication markets experienced a distinctly steeper increase in interest
expense after adoption, with the gap opening immediately at $k=0$ and widening
persistently through $k=+9$.  The triple-interaction coefficient is $+0.044$
percentage points per quarter ($p = 0.020$), implying that high-CECL banks in
high-sophistication markets paid an additional 17--18 basis points annually in
interest expense on top of the average funding cost increase documented in the
main results.  Sophisticated depositors do not simply withdraw---they extract a
larger rate premium as the price of staying, consistent with the dual quantity-and-price
mechanism that \citet{egan2017deposit} document in deposit competition: when risk
is revealed, the marginal depositor raises her required rate rather than
immediately withdrawing, and the bank must pay up or lose the funding.

Panel~D confirms that this additional cost pressure translated into meaningfully
larger profitability losses.  The triple-interaction coefficient for return on
assets is $-0.076$ percentage points per quarter ($p = 0.010$), an annualized
drag of approximately 30 basis points beyond the average CECL effect.  For
community banks whose median ROA runs between 100 and 120 basis points annually,
this represents a substantial incremental earnings penalty concentrated at
institutions whose depositor bases are most capable of interpreting and responding
to credit risk disclosures.  Net interest margin follows a similar pattern
($-0.031$, $p = 0.175$, Panel~E), with the visual separation between high- and
low-sophistication banks widening materially from $k=+3$ onward even if the
full-period average is imprecise.

\paragraph{Interpretation.}
Taken together, the sophistication heterogeneity results support a two-channel
discipline mechanism.  Sophisticated depositors respond to the CECL disclosure
through both quantity---gradually reducing uninsured balances as time deposits
mature---and price, immediately demanding higher rates as compensation for
remaining.  The price channel is statistically sharper because it operates
continuously from the moment of disclosure, whereas the quantity channel is
constrained by the term structure of existing CD contracts and therefore builds
over multiple quarters before becoming fully apparent in average effects.  The
combination of higher funding costs and eventual deposit outflows at
high-sophistication banks amplifies the financial penalties facing high-CECL
institutions beyond what the pooled estimates capture, reinforcing the
interpretation that what we are measuring is active market discipline by informed
creditors rather than a passive mechanical response to the CECL accounting shock.

\subsection{Summary} [@@@ delete this section]
\label{sec:results:summary}

Taken together, the evidence in
Tables~\ref{tab:small_bank_did_time_deposits}--\ref{tab:small_bank_did_financial_outcomes}
and Figures~\ref{fig:es_time_deps_cecl}--\ref{fig:es_financial_outcomes_cecl}
supports a coherent depositor discipline narrative triggered by the CECL
information shock.  At the disclosure date, high-exposure community banks revealed
that their loan portfolios carried materially larger unrecognized expected credit
losses than outsiders had previously estimated.  Uninsured depositors responded
by reducing their balances, exploiting the liquidity of maturing time deposits
to act on the new information.  High-CECL banks replaced the outflow by bidding
up rates on insured time deposits---the only funding tier not subject to run risk
once risk is revealed---paying an enduring premium in funding costs that
compressed both net interest margins and return on assets.

The triple-difference specification with bank $\times$ quarter fixed effects
confirms that this is a within-bank, within-quarter composition shift in the
structure of the deposit portfolio rather than a uniform funding retreat: by
absorbing all bank-specific time-varying unobservables, it rules out any story in
which CECL simply depressed aggregate deposit demand at high-exposure
institutions---the confound that \citet{covitz2004market} identify as the central
threat to discipline tests based on single-outcome regressions.  Parallel
pre-trends in all event-study plots support a causal interpretation.

The results establish that uninsured depositor discipline operates in the
non-crisis, non-TBTF setting where governance theory predicts it should be
strongest \citep{calomiris1991role, flannery1996evidence}.  Prior evidence for the
mechanism is concentrated in crisis episodes \citep{martinez2001depositors,
iyer2016tale} or attenuated by implicit government guarantees \citep{bennett2015market,
berger2015depositors}.  Community banks, which carry no TBTF subsidy
\citep{jacewitz2018deposit, o1990deposit}, face the disciplinary mechanism at
full force---and the CECL disclosure is sharp enough to activate it even absent
acute financial stress.
